
\section{AT-04: Hyperfine structure, Land\'e factors and fields}
\source{2025 PHYS4055 Exam, Question 1(a)--(e); Atomic Assignment 3, Question 4.}
\subsection*{Step 1: find the hyperfine multiplets}
For electronic term \({}^2P_{1/2}\), \(J=1/2\);
let nuclear spin be \(I=5/2\). Coupling
\(\mathbf F=\mathbf I+\mathbf J\) yields
\[
\boxed{F=|I-J|,\ldots,I+J=2,3}.
\]
Using \(F^2=I^2+J^2+2\mathbf I\cdot\mathbf J\)
and \(H_{\rm hfs}=A\mathbf I\cdot\mathbf J/\hbar^2\),
\[
\boxed{E_F=\frac A2\{F(F+1)-I(I+1)-J(J+1)\}}.
\]
Since \(I(I+1)=35/4,\ J(J+1)=3/4\),
\[
\boxed{E_2=-\frac{7A}{4},\quad
E_3=\frac{5A}{4},\quad E_3-E_2=3A}.
\]
For \(A>0\), \(F=3\) lies above \(F=2\); for
\(A<0\), the order reverses.

\subsection*{Step 2: derive and evaluate \(g_J\)}
With \(g_L=1,\ g_S\simeq2\) and \(\mathbf J=\mathbf L+\mathbf S\),
projection of the electronic moment onto \(J\) gives
\[
g_J=1+\frac{J(J+1)+S(S+1)-L(L+1)}{2J(J+1)}.
\]
For \({}^2P_{1/2}\), \(L=1,\ S=1/2,\ J=1/2\):
\[
\boxed{g_J=1+\frac{\frac34+\frac34-2}{2(\frac34)}
=\frac23}.
\]

\subsection*{Step 3: weak-field splitting}
Ignoring the smaller nuclear magnetic moment,
\[
g_F=g_J\frac{F(F+1)+J(J+1)-I(I+1)}{2F(F+1)}
\]
for \(F>0\). Substitution gives
\[
\boxed{g_{F=2}=-\frac19,\qquad g_{F=3}=+\frac19.}
\]
At weak \(B\), the energies are
\(E_{F,m_F}\simeq E_F+\mu_Bg_Fm_FB\).
Thus \(F=2\) splits into five magnetic substates
and \(F=3\) into seven; for generic weak nonzero
\(B\) the 12 states have different energies.
The slopes of the two manifolds have opposite signs.

\subsection*{Step 4: describe a strong field}
When Zeeman energy dominates hyperfine coupling,
the \(\ket{F,m_F}\) basis is no longer approximately
diagonal. Use \(\ket{Jm_J,Im_I}\) instead:
\[
E\simeq \mu_Bg_Jm_JB+A\,m_Im_J
\quad\text{(neglecting nuclear Zeeman energy)}.
\]
There are two \(m_J=\pm1/2\) electronic branches,
each with six \(m_I\) values.
The electric-quadrupole hyperfine operator is rank 2;
the electronic \(J=1/2\to J=1/2\) reduced matrix
element vanishes because \(J+J=1<2\) (triangle rule).
\checkmethod{A strong-field calculation must switch coupling
basis, not simply extend the weak-field straight lines
indefinitely. The \(2F+1\) degeneracies must sum to
\((2I+1)(2J+1)=12\).}
